% GENERATED FILE. Edit canonical lessons, then run npm run generate:books.
% canonical-source-hash: fnv1a64:ecf8b8b2779be64d
% canonical-lessons: GU-C409-paronagat, GU-C409-menejar, GU-C409-sagvad, GU-C409-agvad, GU-C409-avkarvu

\chapter{Lodgings}
\label{ch:gu-lodgings}
\hlchaptermodality{409}

\begin{chapteropening}
\textbf{By the end of this chapter:} \emph{I can talk about what a lodging offers.}

Complete the last lesson of chapter 409: I can talk about what a lodging offers.
\end{chapteropening}

\section[\texorpdfstring{\gu{પરોણાગત} (paroṇāgat)}{પરોણાગત (paroṇāgat)}]{\gu{પરોણાગત} (paroṇāgat) — hospitality (to a guest)}

\label{lesson:GU-C409-paronagat}



\begin{quote}
Before the new one: say the Gujarati for a rest on the way, then the Gujarati for to walk up and down.
\end{quote}

\subsection*{You'll want to know: \gu{પરોણાગત}}
\textbf{\gu{પરોણાગત}} — \emph{paroṇāgat} — \textquotedblleft{}hospitality (to a guest)\textquotedblright{}.

\begin{quote}
\gu{આ} \gu{હોટેલની} \gu{પરોણાગત} \gu{ઘણી} \gu{સારી} \gu{છે}.
\end{quote}

\begin{quote}
\emph{ā hoṭelnī paroṇāgat ghaṇī sārī che.}
\end{quote}

\begin{quote}
The hospitality at this hotel is very good.
\end{quote}

A feminine noun.

\begin{cousinweb}
\textbf{\gu{પરોણાગત}} (\emph{paroṇāgat}) is built on \gu{પરોણો}, a guest, which CDIAL lists under Sanskrit \emph{prāhuṇa}, \textquotedblleft{}a guest\textquotedblright{} (CDIAL 8973): the care a house gives to the guest who stays in it.
\end{cousinweb}

\begin{grammarlens}[title={What you've built}]
One more word for getting around.
\end{grammarlens}

\subsection*{Your turn}
\begin{itemize}
\raggedright
  \item \emph{Say it:} \emph{paroṇāgat}
  \item \emph{Say it:} \emph{paroṇāgat}, once more
  \item \emph{Say it:} \emph{ā hoṭelnī paroṇāgat ghaṇī sārī che.} — The hospitality at this hotel is very good.
  \item \emph{Recall:} say the Gujarati for a rest on the way, then the Gujarati for to walk up and down, then say \emph{paroṇāgat} again
\end{itemize}

\begin{tcolorbox}[breakable,colback=teal!4,colframe=teal!35!black,title={Before you move on}]
What does \gu{પરોણાગત} mean? (\textquotedblleft{}Hospitality (to a guest)\textquotedblright{}.) Say it once more.
\end{tcolorbox}
\section[\texorpdfstring{\gu{મેનેજર} (menejar)}{મેનેજર (menejar)}]{\gu{મેનેજર} (menejar) — a manager}

\label{lesson:GU-C409-menejar}



\begin{quote}
Before the new one: say the Gujarati for to walk up and down, then the Gujarati for hospitality.
\end{quote}

\subsection*{You'll want to know: \gu{મેનેજર}}
\textbf{\gu{મેનેજર}} — \emph{menejar} — \textquotedblleft{}a manager\textquotedblright{}.

\begin{quote}
\gu{હોટેલના} \gu{મેનેજરે} \gu{અમને} \gu{ચાવી} \gu{આપી}.
\end{quote}

\begin{quote}
\emph{hoṭelnā menejre amne cāvī āpī.}
\end{quote}

\begin{quote}
The hotel manager gave us the key.
\end{quote}

A masculine noun.

\begin{cousinweb}
\textbf{\gu{મેનેજર}} (\emph{menejar}) is borrowed from English \emph{manager}.
\end{cousinweb}

\begin{grammarlens}[title={What you've built}]
One more word for getting around.
\end{grammarlens}

\subsection*{Your turn}
\begin{itemize}
\raggedright
  \item \emph{Say it:} \emph{menejar}
  \item \emph{Say it:} \emph{menejar}, once more
  \item \emph{Say it:} \emph{hoṭelnā menejre amne cāvī āpī.} — The hotel manager gave us the key.
  \item \emph{Recall:} say the Gujarati for to walk up and down, then the Gujarati for hospitality, then say \emph{menejar} again
\end{itemize}

\begin{tcolorbox}[breakable,colback=teal!4,colframe=teal!35!black,title={Before you move on}]
What does \gu{મેનેજર} mean? (\textquotedblleft{}A manager\textquotedblright{}.) Say it once more.
\end{tcolorbox}
\section[\texorpdfstring{\gu{સગવડ} (sagvaḍ)}{સગવડ (sagvaḍ)}]{\gu{સગવડ} (sagvaḍ) — a convenience; a facility}

\label{lesson:GU-C409-sagvad}



\begin{quote}
Before the new one: say the Gujarati for hospitality, then the Gujarati for a manager.
\end{quote}

\subsection*{You'll want to know: \gu{સગવડ}}
\textbf{\gu{સગવડ}} — \emph{sagvaḍ} — \textquotedblleft{}a convenience; a facility\textquotedblright{}.

\begin{quote}
\gu{આ} \gu{હોટેલમાં} \gu{બધી} \gu{સગવડ} \gu{છે}.
\end{quote}

\begin{quote}
\emph{ā hoṭelmā̃ badhī sagvaḍ che.}
\end{quote}

\begin{quote}
This hotel has every convenience.
\end{quote}

A feminine noun.

\begin{cousinweb}
\textbf{\gu{સગવડ}} (\emph{sagvaḍ}) No dictionary we use records an older form, so its history is not settled.
\end{cousinweb}

\begin{grammarlens}[title={What you've built}]
One more word for getting around.
\end{grammarlens}

\subsection*{Your turn}
\begin{itemize}
\raggedright
  \item \emph{Say it:} \emph{sagvaḍ}
  \item \emph{Say it:} \emph{sagvaḍ}, once more
  \item \emph{Say it:} \emph{ā hoṭelmā̃ badhī sagvaḍ che.} — This hotel has every convenience.
  \item \emph{Recall:} say the Gujarati for hospitality, then the Gujarati for a manager, then say \emph{sagvaḍ} again
\end{itemize}

\begin{tcolorbox}[breakable,colback=teal!4,colframe=teal!35!black,title={Before you move on}]
What does \gu{સગવડ} mean? (\textquotedblleft{}A convenience; a facility\textquotedblright{}.) Say it once more.
\end{tcolorbox}
\section[\texorpdfstring{\gu{અગવડ} (agvaḍ)}{અગવડ (agvaḍ)}]{\gu{અગવડ} (agvaḍ) — an inconvenience; a discomfort}

\label{lesson:GU-C409-agvad}



\begin{quote}
Before the new one: say the Gujarati for a manager, then the Gujarati for a convenience.
\end{quote}

\subsection*{You'll want to know: \gu{અગવડ}}
\textbf{\gu{અગવડ}} — \emph{agvaḍ} — \textquotedblleft{}an inconvenience; a discomfort\textquotedblright{}.

\begin{quote}
\gu{નાના} \gu{ઓરડામાં} \gu{અમને} \gu{ઘણી} \gu{અગવડ} \gu{પડી}.
\end{quote}

\begin{quote}
\emph{nānā orḍāmā̃ amne ghaṇī agvaḍ paḍī.}
\end{quote}

\begin{quote}
We were very uncomfortable in the small room.
\end{quote}

A feminine noun.

\begin{cousinweb}
\textbf{\gu{અગવડ}} (\emph{agvaḍ}) is the opposite of \gu{સગવડ}, a convenience, and the two are usually heard together. No dictionary we use records an older form, so its history is not settled.
\end{cousinweb}

\begin{grammarlens}[title={What you've built}]
One more word for getting around.
\end{grammarlens}

\subsection*{Your turn}
\begin{itemize}
\raggedright
  \item \emph{Say it:} \emph{agvaḍ}
  \item \emph{Say it:} \emph{agvaḍ}, once more
  \item \emph{Say it:} \emph{nānā orḍāmā̃ amne ghaṇī agvaḍ paḍī.} — We were very uncomfortable in the small room.
  \item \emph{Recall:} say the Gujarati for a manager, then the Gujarati for a convenience, then say \emph{agvaḍ} again
\end{itemize}

\begin{tcolorbox}[breakable,colback=teal!4,colframe=teal!35!black,title={Before you move on}]
What does \gu{અગવડ} mean? (\textquotedblleft{}An inconvenience; a discomfort\textquotedblright{}.) Say it once more.
\end{tcolorbox}
\section[\texorpdfstring{\gu{આવકારવું} (āvkārvũ)}{આવકારવું (āvkārvũ)}]{\gu{આવકારવું} (āvkārvũ) — to welcome (a guest)}

\label{lesson:GU-C409-avkarvu}



\begin{quote}
Before the new one: say the Gujarati for a convenience, then the Gujarati for an inconvenience.
\end{quote}

\subsection*{You'll want to know: \gu{આવકારવું}}
\textbf{\gu{આવકારવું}} — \emph{āvkārvũ} — \textquotedblleft{}to welcome (a guest)\textquotedblright{}.

\begin{quote}
\gu{મેનેજર} \gu{દરેક} \gu{મહેમાનને} \gu{હસીને} \gu{આવકારે} \gu{છે}.
\end{quote}

\begin{quote}
\emph{menejar darek mahemānne hasīne āvkāre che.}
\end{quote}

\begin{quote}
The manager welcomes every guest with a smile.
\end{quote}

Past, with \gu{તેણે} (he or she …): \gu{તેણે} \gu{મહેમાનને} \gu{આવકાર્યો}, \emph{teṇe mahemānne āvkāryo}.

\begin{cousinweb}
\textbf{\gu{આવકારવું}} (\emph{āvkārvũ}) is built on \gu{આવકાર}, \textquotedblleft{}a welcome\textquotedblright{}, which seems to go with \gu{આવવું}, to come. No dictionary we use records an older form, so its history is not settled.
\end{cousinweb}

\begin{grammarlens}[title={What you've built}]
One more word for getting around.
\end{grammarlens}

\subsection*{Your turn}
\begin{itemize}
\raggedright
  \item \emph{Say it:} \emph{āvkārvũ}
  \item \emph{Say it:} \emph{āvkārvũ}, once more
  \item \emph{Say it:} \emph{menejar darek mahemānne hasīne āvkāre che.} — The manager welcomes every guest with a smile.
  \item \emph{Recall:} say the Gujarati for a convenience, then the Gujarati for an inconvenience, then say \emph{āvkārvũ} again
\end{itemize}

\begin{tcolorbox}[breakable,colback=teal!4,colframe=teal!35!black,title={Before you move on}]
What does \gu{આવકારવું} mean? (\textquotedblleft{}To welcome (a guest)\textquotedblright{}.) Say it once more.
\end{tcolorbox}
